\section{NP-completeness of the Min Cut-Path Problem}
\label{sec:np-completeness}
We prove that the decision version of the \MinCutPath{} problem is NP-complete by a reduction from the \ThreeSAT{} problem.

\begin{problem}[\ThreeSAT]
Let $U = \{x_1, x_2, \dots, x_n\}$ be a finite set of Boolean variables, and let
$\mathcal{C} = \{C_1, C_2, \dots, C_m\}$ be a collection of clauses, where each clause
$C_k$ is a disjunction of exactly three literals, and each literal is either a variable $x_i$ or its negation $\neg x_i$.

The \ThreeSAT{} decision problem is defined as follows:

\begin{center}
\begin{tabular}{p{0.12\linewidth}p{0.8\linewidth}}
\textbf{Input:} & A finite set of Boolean variables $U$ and a collection $\mathcal{C}$ of clauses, each containing exactly three literals.\\
\textbf{Question:} & Does there exist a truth assignment $\tau: U \rightarrow \{\mathsf{true}, \mathsf{false}\}$ such that every clause $C_k \in \mathcal{C}$ evaluates to \textsf{true} under $\tau$?
\end{tabular}
\end{center}
\end{problem}

\ThreeSAT{} is NP-complete~\cite{Cook1971Complexity}.
Throughout this section, $n$ and $m$ denote the numbers of variables and clauses of the \ThreeSAT{} instance under consideration.

The proof has two steps and uses an intermediate problem, \SSP.
First, we show that \SSP{} is NP-complete by a polynomial-time reduction from \ThreeSAT{} (Section~\ref{sec:ssp}).
Second, we reduce \SSP{} to \MinCutPath{} (Section~\ref{sec:ssp-to-mcp}).

\subsection{The Separating Shortest Path Problem}\label{sec:ssp}
The \SSP{} problem asks whether the removal of some shortest $u$--$v$ path disconnects $u$ from $v$.

\begin{problem}[\SSP]
Let $G = (V,E)$ be an undirected graph and let $u,v \in V$ be two distinct vertices.
A shortest $u$--$v$ path $P$ is called a \emph{separating shortest path} if the graph $G \setminus P$ contains no $u$--$v$ path.
The \SSP{} decision problem is defined as follows:

\begin{center}
\begin{tabular}{p{0.12\linewidth}p{0.8\linewidth}}
\textbf{Input:} & A graph $G = (V,E)$ and two distinct vertices $u,v \in V$.\\
\textbf{Question:} & Does there exist a shortest $u$--$v$ path in $G$ that is a separating shortest path?
\end{tabular}
\end{center}
\end{problem}

Unlike \MinCutPath, \SSP{} is not an optimization problem: a solution is a single shortest $u$--$v$ path, that is, a cut-path of size $d(u,v)$ (see the proof of Theorem~\ref{thm:mcp-np-complete}).
This makes the reduction from \ThreeSAT{} more direct.

\subsubsection{Gadgets and Basic Operations for the Reduction}

The reduction from \ThreeSAT{} to \SSP{} is built from \emph{gadgets} and \emph{basic operations} on them.
The construction consists of two parts:
\begin{enumerate}
    \item the \emph{base structure}, also referred to as the \emph{chain};
    \item the \emph{threads}, which are threaded through the base structure.
\end{enumerate}
Figure~\ref{fig:chain-threads} shows both parts.

\begin{figure}[H]
    \centering
    \includegraphics{chain-threads.pdf}
    \caption{Layout of the construction: the chain (black) and a thread (red)}
    \label{fig:chain-threads}
\end{figure}

\paragraph{Chain.}
The smallest gadget is the \emph{chain link}, a cycle with two distinguished vertices $p$ and $q$ that split it into two paths of equal length.

\begin{definition}[Chain Link]
\label{def:chain-link}
Let $G=(V,E)$ be a graph and let $p,q \in V$ be two distinct vertices.
A \emph{$(p,q)$-chain link} is a subgraph $L \subseteq G$ such that:
\begin{enumerate}
    \item $L$ is a simple cycle in $G$;
    \item the vertices $p$ and $q$ lie on $L$ and partition it into two internally disjoint $p$--$q$ paths;
    \item these two $p$--$q$ paths have equal length.
\end{enumerate}
We refer to one of the two $p$--$q$ paths as the \emph{positive path} of the chain link and to the other as the \emph{negative path}.
When a chain link is used inside a chain, its boundary vertices are exactly $p$ and~$q$, i.e., the points where the adjacent single edges of the chain are attached.
\end{definition}

Figure~\ref{fig:chain-link} shows a chain link.

\begin{figure}[H]
    \centering
    \includegraphics{chain-link.pdf}
    \caption{A chain link}
    \label{fig:chain-link}
\end{figure}

A chain alternates single edges and chain links; it begins and ends with a single edge.

\begin{definition}[Chain]
\label{def:chain}
Let $G=(V,E)$ be a graph, and let $u,v\in V$.
We say that $G$ is a \emph{$u$--$v$ chain} with $r$ chain links if it is the union of edge-disjoint subgraphs
\[
G = H_1 \cup L_1 \cup H_2 \cup L_2 \cup \cdots \cup L_r \cup H_{r+1},
\]
where:
\begin{enumerate}
    \item each $H_i$ consists of a single edge with end-vertices $p_i,q_i$;
    \item each $L_i$ is a $(q_i,p_{i+1})$-chain link;
    \item the first end-vertex $p_1$ coincides with $u$ and the last end-vertex $q_{r+1}$ coincides with $v$;
    \item two of these subgraphs share a vertex only if they are consecutive, and then they share exactly one vertex ($q_i$ or $p_{i+1}$).
\end{enumerate}
Every $u$--$v$ path in a chain traverses all single edges and, in each chain link, either its positive or its negative path.
\end{definition}

Figure~\ref{fig:chain} shows a chain.

\begin{figure}[H]
    \centering
    \includegraphics{chain.pdf}
    \caption{A chain}
    \label{fig:chain}
\end{figure}

Let $C_1, C_2, \dots, C_m$ denote the clauses of the input formula and let $x_1, x_2, \dots, x_n$ denote its variables.
In the reduction, every chain link has one of three types:
\begin{enumerate}
    \item \textbf{Initialization chain link.}
    For each variable $x_i$ we introduce an initialization chain link, denoted $I_i$.
    This link marks the beginning of the part of the chain that corresponds to the variable $x_i$.
    The total number of initialization chain links is $n$.

    \item \textbf{Terminal chain link.}
    For each variable $x_i$ we also introduce a terminal chain link, denoted $T_i$.
    This link marks the end of the part of the chain associated with the variable $x_i$.
    The total number of terminal chain links is $n$.

    \item \textbf{Literal chain link.}
    For each clause $C_k = (\ell_1 \vee \ell_2 \vee \ell_3)$ and each $j \in \{1,2,3\}$, let $x_i$ be the variable of the literal $\ell_j$.
    We introduce a literal chain link, denoted $L_{j,k,i}$ (or $L_{j,k}$ for short).
    This link represents the occurrence of the variable $x_i$ as the $j$-th literal of the clause $C_k$.
    The total number of literal chain links is $3m$.
\end{enumerate}
In total, the construction contains $3m + 2n$ chain links.
The chain links that belong to one variable are consecutive in the chain, as illustrated in Figure~\ref{fig:chain-link-types}.

\begin{figure}[H]
    \centering
    \includegraphics{chain-link-types.pdf}
    \caption{Chain links of a variable $x_i$: the initialization link, the literal links, and the terminal link}
    \label{fig:chain-link-types}
\end{figure}

\paragraph{Threads and threading.}
\emph{Threads} are paths that pass through the chain links; inserting such a path into the chain is called \emph{threading}.

\begin{definition}[Thread]
\label{def:thread}
Let $G=(V,E)$ be a graph containing a $u$--$v$ chain with chain links $L_1,L_2,\dots,L_r$.
A \emph{thread} is a $u$--$v$ path $P$ in $G$ such that:
\begin{enumerate}
    \item for each chain link $L_i$, the path $P$ shares at most one edge with $L_i$, and $P$ contains none of the single edges of the chain;
    \item $P$ is edge-disjoint from every other thread and shares only the vertices $u$ and $v$ with it.
\end{enumerate}
\end{definition}

Threading subdivides an edge of a chain link that no thread uses yet into three edges and routes the thread through the middle one.

\begin{definition}[Threading]
\label{def:threading}
Let $L$ be a chain link in a $u$--$v$ chain, and let $P$ be a thread not yet passing through $L$.
The \emph{threading operation} on $(P,L)$ consists of the following elementary steps:
\begin{enumerate}
    \item Select an edge $e=\{x,y\}$ of $L$ that is not used by any existing thread.
    \item Subdivide $e$ into a path $x$--$z_1$--$z_2$--$y$ of three edges, introducing two new vertices $z_1,z_2$.
    \item Extend $P$ so that it traverses the middle edge $\{z_1,z_2\}$, called the \emph{crossing edge}, as its unique crossing of the chain link $L$.
\end{enumerate}
If the selected edge $e$ belongs to the positive path of $L$, we say that $P$ is \emph{positively threaded through $L$}; if $e$ belongs to the negative path of $L$, we say that $P$ is \emph{negatively threaded through $L$}.
In general, we say that the thread $P$ is \emph{threaded through $L$ at the edge $e$}.
\end{definition}

Figure~\ref{fig:threading} shows the threading operation.

\begin{figure}[H]
    \centering
    \includegraphics{threading.pdf}
    \caption{The threading operation}
    \label{fig:threading}
\end{figure}

Every thread has one of three types:
\begin{enumerate}
    \item \textbf{2-synchronization thread.}
    This thread passes through an initialization chain link $I_i$ and the corresponding terminal chain link $T_i$ (Figure~\ref{fig:two-synchronization}).
    Threads of this type force every separating shortest path to traverse paths of the same sign in $I_i$ and $T_i$ (Lemmas~\ref{lem:chain-paths} and~\ref{lem:synchronization}).

    \item \textbf{3-synchronization thread.}
    This thread passes through the initialization chain link $I_i$, a literal chain link $L_{j,k,i}$, and the terminal chain link $T_i$ corresponding to the same variable $x_i$ (Figure~\ref{fig:three-synchronization}).
    Threads of this type force every separating shortest path to traverse the path of this sign in $L_{j,k,i}$ as well (Lemma~\ref{lem:synchronization}).

    \item \textbf{Clause thread.}
    For each clause $C_k = (\ell_{1} \vee \ell_{2} \vee \ell_{3})$, a clause thread is introduced that passes through the three literal chain links $L_{1,k}$, $L_{2,k}$, $L_{3,k}$ corresponding to $\ell_{1}, \ell_{2}, \ell_{3}$ (Figure~\ref{fig:clause-thread}).
    It encodes the clause: a separating shortest path exists only if at least one of the three literals is satisfied (Lemma~\ref{lem:clause-threads}).
\end{enumerate}

\begin{figure}[H]
    \centering
    \includegraphics{two-synchronization-threads.pdf}
    \caption{The two 2-synchronization threads of a variable $x_i$}
    \label{fig:two-synchronization}
\end{figure}

\begin{figure}[H]
    \centering
    \includegraphics{three-synchronization-threads.pdf}
    \caption{The two 3-synchronization threads of a literal chain link $L_{j,k,i}$}
    \label{fig:three-synchronization}
\end{figure}

\begin{figure}[H]
    \centering
    \includegraphics{clause-thread.pdf}
    \caption{The clause thread of a clause $C_k = (x_1 \vee \neg x_2 \vee x_3)$}
    \label{fig:clause-thread}
\end{figure}

\subsubsection{The Construction}
\label{sec:reduction-construction}
We prove that \SSP{} is NP-complete (Theorem~\ref{thm:ssp-np-complete}) by a polynomial-time reduction from \ThreeSAT{} that uses the chain and the threads.
Given an instance of \ThreeSAT, we construct a graph $G$ with two distinguished vertices $u,v$ such that $G$ contains a separating shortest $u$--$v$ path if and only if the given formula is satisfiable.

\paragraph{Auxiliary procedures.}
The reduction uses the following auxiliary procedures:
\begin{itemize}
    \item $\textsc{BuildChain}(u,v,r)$: constructs a $u$--$v$ chain with $r$ chain links, where each chain link is realized as a $4$-cycle.
    \item $\textsc{Thread}(u,v,[(K_1,\sigma_1),\dots,(K_t,\sigma_t)])$: creates a new thread from $u$ to $v$ that passes through the chain links $K_1,\dots,K_t$ in this order, where $\sigma_i \in \{+,-\}$ specifies whether the thread is threaded through $K_i$ positively or negatively.
    The $t$ crossing edges are joined to each other and to the vertices $u$ and $v$ by $t+1$ new, internally vertex-disjoint paths, which we call the \emph{connecting paths} of the thread.
    \item $\textsc{Calibrate}(G)$: equalizes the lengths of the relevant paths in two steps.
    First, in every chain link whose positive and negative paths have different lengths, an edge of the shorter path that is not a crossing edge is subdivided repeatedly until both paths have the same length.
    Second, let $\Lambda$ denote the number of edges of a $u$--$v$ path in the resulting chain; every connecting path of every thread is subdivided until it has at least $\Lambda$ edges.
\end{itemize}

\paragraph{Auxiliary notation.}
We also introduce the following notational conventions:
\begin{itemize}
    \item Occurrences: for each variable $x_i$, we denote by $O_i$ the set of all pairs $(j,k)$ such that the $j$-th literal of the clause $C_k$ is $x_i$ or $\neg x_i$.
    \item Sign function: for each literal $\ell$ we define
\[
s(\ell) =
\begin{cases}
+ & \text{if $\ell = x_i$ for some variable $x_i$,}\\
- & \text{if $\ell = \neg x_i$ for some variable $x_i$.}
\end{cases}
\]
\end{itemize}

\paragraph{The graph of the reduction.}
Algorithm~\ref{alg:reduction} constructs the graph $G$ and the vertices $u,v$.
From now on, $G$, $u$ and $v$ denote its output.

\begin{algorithm}[H]
\caption{Reduction from \ThreeSAT{} to \SSP}
\label{alg:reduction}
\begin{algorithmic}[1]
\Require Variables $U=\{x_1,\dots,x_n\}$, clauses $\mathcal{C}=\{C_1,\dots,C_m\}$
\Ensure Graph $G$, vertices $u,v$
\State $G \gets \textsc{BuildChain}(u,v,2n+3m)$
\For{$i \gets 1$ to $n$}
    \State reserve the next $2+|O_i|$ free chain links
    \State label the first of them as $I_i$ and the last of them as $T_i$
    \For{each $(j,k) \in O_i$}
        \State label the next free reserved link as $L_{j,k,i}$
    \EndFor
\EndFor
\For{$i \gets 1$ to $n$}
    \State $\textsc{Thread}(u,v,[(I_i,+),(T_i,-)])$
    \State $\textsc{Thread}(u,v,[(I_i,-),(T_i,+)])$
    \For{each $(j,k) \in O_i$}
        \State $\textsc{Thread}(u,v,[(I_i,+),(L_{j,k,i},-),(T_i,+)])$
        \State $\textsc{Thread}(u,v,[(I_i,-),(L_{j,k,i},+),(T_i,-)])$
    \EndFor
\EndFor
\For{$k \gets 1$ to $m$}
    \State $\ell_1,\ell_2,\ell_3 \gets$ the literals of $C_k$
    \State $\textsc{Thread}(u,v,[(L_{1,k},s(\ell_1)),(L_{2,k},s(\ell_2)),(L_{3,k},s(\ell_3))])$
\EndFor
\State $\textsc{Calibrate}(G)$
\State \Return $G$, $u$, $v$
\end{algorithmic}
\end{algorithm}

In words, the chain consists of $n$ consecutive blocks, one for each variable; the block of $x_i$ starts with $I_i$, ends with $T_i$, and contains the literal chain links of all occurrences of $x_i$ in between.
Each variable receives a pair of 2-synchronization threads, each occurrence of a variable receives a pair of 3-synchronization threads, and each clause receives one clause thread.
Every thread passes through at least two chain links, so it has at least two crossing edges.

\paragraph{Lengths after the calibration.}
The final calibration step guarantees that both paths of every chain link have the same length and that the connecting paths of the threads are sufficiently long.
Hence the subdivided $4$-cycles are again chain links, the chain in $G$ is a $u$--$v$ chain in the sense of Definition~\ref{def:chain}, every $u$--$v$ path of this chain has exactly $\Lambda$ edges, and every connecting path has at least $\Lambda$ edges.

\subsubsection{Correctness of the Reduction}
\label{sec:reduction-correctness}
We call the $u$--$v$ paths of the chain in $G$ \emph{chain paths}; a chain path is determined by choosing, in every chain link, either its positive or its negative path.
We say that a chain path $P$ \emph{hits} a thread if $P$ and the thread share an edge.
Since a thread shares only its crossing edges with the chain, a chain path $P$ hits a thread created by $\textsc{Thread}(u,v,[(K_1,\sigma_1),\dots,(K_t,\sigma_t)])$ if and only if, for some $i$, the path $P$ traverses the path of $K_i$ with the sign $\sigma_i$.

Every thread is a $u$--$v$ path in $G$.
Therefore, a chain path $P$ can be separating only if it hits every thread.

A chain path is \emph{consistent} if, for every variable $x_i$, it traverses paths of the same sign in all chain links of the block of $x_i$.
Consistent chain paths correspond to truth assignments: we set $\tau(x_i)=\mathsf{true}$ if the path traverses the positive paths in the block of $x_i$, and $\tau(x_i)=\mathsf{false}$ otherwise.
Conversely, every truth assignment $\tau$ determines a consistent chain path, which traverses, in all chain links of the block of $x_i$, the positive path if $\tau(x_i)=\mathsf{true}$ and the negative path if $\tau(x_i)=\mathsf{false}$.
We call it the \emph{consistent chain path of~$\tau$}.

\begin{lemma}[Shortest paths]
\label{lem:chain-paths}
All chain paths have $\Lambda$ edges, and every other $u$--$v$ path in $G$ has more than $\Lambda$ edges.
In particular, the shortest $u$--$v$ paths in $G$ are exactly the chain paths.
\end{lemma}

\begin{proof}
After the calibration, all chain paths have the same length $\Lambda$.
Let $Q$ be a $u$--$v$ path in $G$ that is not a chain path.
By Definition~\ref{def:chain}, every $u$--$v$ path that uses only edges of the chain is a chain path, and the edges of $G$ outside the chain are the edges of the connecting paths.
Hence $Q$ contains an edge of some connecting path.
The internal vertices of a connecting path have degree two in $G$, because the procedure $\textsc{Thread}$ creates the connecting paths as new, internally vertex-disjoint paths and $\textsc{Calibrate}$ subdivides them only by new vertices of degree two.
Since $u$ and $v$ are not internal vertices of a connecting path, $Q$ contains the whole connecting path, which has at least $\Lambda$ edges.
Every thread has at least two crossing edges, so every connecting path has an end-vertex on a crossing edge, and no connecting path joins $u$ and $v$ directly.
Hence $Q$ contains at least one further edge, and $|Q| > \Lambda$.
\end{proof}

\begin{lemma}[Synchronization]
\label{lem:synchronization}
A chain path hits all synchronization threads if and only if it is consistent.
\end{lemma}

\begin{proof}
Let $P$ be a chain path and let $x_i$ be a variable.
By Algorithm~\ref{alg:reduction}, the block of $x_i$ consists of $I_i$, the literal chain links $L_{j,k,i}$ with $(j,k) \in O_i$, and $T_i$.
\begin{itemize}
    \item \emph{2-synchronization threads:} the chain path $P$ hits both 2-synchronization threads of the variable $x_i$ (Figure~\ref{fig:two-synchronization}) if and only if it traverses paths of the same sign in $I_i$ and in $T_i$.
    Indeed, the first thread requires the positive path in $I_i$ or the negative path in $T_i$, and the second thread requires the negative path in $I_i$ or the positive path in $T_i$.
    \item \emph{3-synchronization threads:} if $P$ traverses paths of the same sign in $I_i$ and $T_i$, then it hits both 3-synchronization threads of a literal link $L_{j,k,i}$ (Figure~\ref{fig:three-synchronization}) if and only if it traverses the path of this sign in $L_{j,k,i}$ as well.
    Indeed, one of the two threads is hit exactly when $P$ uses the positive path in $I_i$ or $T_i$ or the negative path in $L_{j,k,i}$, and the other one exactly when $P$ uses the negative path in $I_i$ or $T_i$ or the positive path in $L_{j,k,i}$.
\end{itemize}
Suppose that $P$ hits all synchronization threads.
By the first item, $P$ traverses paths of the same sign in $I_i$ and $T_i$; by the second item, it traverses the path of this sign in every $L_{j,k,i}$ as well.
Conversely, if $P$ traverses paths of the same sign in all chain links of the block of $x_i$, then both items show that $P$ hits all synchronization threads of $x_i$.
Consequently, a chain path hits all synchronization threads if and only if, for every variable $x_i$, it traverses paths of the same sign in all chain links of the block of $x_i$, that is, if and only if it is consistent.
\end{proof}

\begin{lemma}[Clause threads]
\label{lem:clause-threads}
Let $P$ be a consistent chain path with the truth assignment $\tau$.
Then $P$ hits all clause threads if and only if $\tau$ satisfies every clause.
\end{lemma}

\begin{proof}
Let $C_k = (\ell_1 \vee \ell_2 \vee \ell_3)$ be a clause.
The path $P$ hits the clause thread of $C_k$ if and only if, for some $j$, it traverses the path of $L_{j,k}$ with the sign $s(\ell_j)$, i.e., if and only if $\tau$ satisfies the literal $\ell_j$.
Indeed, if $x_i$ is the variable of $\ell_j$, then $P$ traverses the positive path of $L_{j,k}$ exactly when $\tau(x_i)=\mathsf{true}$, and $s(\ell_j) = +$ exactly when $\ell_j = x_i$.
Thus a consistent chain path hits all clause threads if and only if its truth assignment satisfies every clause.
\end{proof}

\begin{lemma}[Separation]
\label{lem:separating}
Let $\tau$ be a truth assignment that satisfies every clause, and let $P$ be the consistent chain path of $\tau$.
Then $P$ is a separating shortest $u$--$v$ path in $G$.
\end{lemma}

\begin{proof}
By Lemma~\ref{lem:chain-paths}, $P$ is a shortest $u$--$v$ path, so we have to show that $G\setminus P$ contains no $u$--$v$ path.

\emph{The graph $G \setminus P$.}
In $G \setminus P$, all single edges of the chain are missing, and in every chain link $K$ only the path not traversed by $P$ remains; we call it the \emph{unused path} of $K$ and denote it by $K^{\circ}$.
A crossing edge that lies on $P$ is removed together with the two edges of the chain link adjacent to it; hence the two connecting paths that end at this crossing edge are \emph{dead ends} in $G \setminus P$.
Indeed, each end-vertex of a crossing edge has degree three in $G$, with two edges in the chain link and one in a connecting path, so in $G \setminus P$ it has degree one.
As it is neither $u$ nor $v$, no $u$--$v$ path of $G \setminus P$ contains it.
A crossing edge on an unused path remains and joins its two adjacent connecting paths to this unused path.
Consequently, a path in $G \setminus P$ can move between connecting paths and unused paths only at crossing edges that are not on $P$.

\emph{The auxiliary multigraph $\Gamma$.}
We record these moves in a multigraph $\Gamma$ whose nodes are $u$, $v$ and the unused paths $K^{\circ}$ of all chain links $K$.
An end of a connecting path is \emph{open} if it is $u$, $v$, or an end-vertex of a crossing edge not on $P$; in the last case it lies on $K^{\circ}$, where $K$ is the chain link of this crossing edge.
Every connecting path with two open ends gives an edge of $\Gamma$ between the nodes of its ends.

If $G \setminus P$ contains a $u$--$v$ path $Q$, then $\Gamma$ contains a $u$--$v$ walk.
Indeed, every edge of $G \setminus P$ lies on an unused path or on a connecting path, because the single edges and the traversed paths of the chain links lie on $P$.
Distinct unused paths are vertex-disjoint and contain neither $u$ nor $v$, because two chain links are never consecutive in the chain (Definition~\ref{def:chain}, item~4) and $u$, $v$ lie only on single edges of the chain.
If $Q$ uses an edge of a connecting path, it contains the whole connecting path, since the internal vertices of a connecting path have degree two (proof of Lemma~\ref{lem:chain-paths}).
Both ends of this connecting path are open, since an end-vertex of a crossing edge on $P$ lies on no $u$--$v$ path of $G \setminus P$.
Hence $Q$ consists of whole connecting paths with open ends and of segments between them, each segment inside a single unused path.
Replacing every such connecting path by its edge of $\Gamma$ and every segment by its node gives a $u$--$v$ walk in $\Gamma$.

\emph{The edges of $\Gamma$.}
Every edge of $G$ belongs to the chain or to a connecting path of a thread, and every thread is a synchronization thread or a clause thread.
The following three cases are therefore exhaustive; they also describe all ways in which a path starting at $u$ could progress in $G \setminus P$.
\begin{enumerate}
  \item \emph{Using a chain edge.}
  The single edge of the chain incident to $u$ lies on $P$, so this option is blocked immediately.
  More generally, all single edges of the chain lie on $P$, and the edges of an unused path $K^{\circ}$ belong to the node $K^{\circ}$ of $\Gamma$; the chain gives no edge of $\Gamma$.

  \item \emph{Using a synchronization thread.}
  Such a thread first reaches $I_i$.
  If it crosses the path of $I_i$ traversed by $P$, it ends there.
  Otherwise, it reaches the unused path of $I_i$.
  The threads that cross this unused path are one 2-synchronization thread, whose next crossing edge is in $T_i$, and one 3-synchronization thread for each literal link $L_{j,k,i}$, whose next crossing edge is in $L_{j,k,i}$; all these next crossing edges lie on $P$.
  Hence every continuation from the unused path of $I_i$ leads back to $u$ or to a dead end.

  To find the edges of $\Gamma$ given by these threads, suppose first that $\tau(x_i)=\mathsf{true}$.
  Then $P$ traverses the positive paths in the block of $x_i$, so a crossing edge in this block lies on $P$ if and only if its thread is threaded positively there.
  The four kinds of synchronization threads of $x_i$ give the following edges of $\Gamma$:
  \begin{itemize}
    \item $[(I_i,+),(T_i,-)]$: only the last connecting path has two open ends; it gives the edge $T_i^{\circ}v$;
    \item $[(I_i,-),(T_i,+)]$: only the first connecting path has two open ends; it gives the edge $uI_i^{\circ}$;
    \item $[(I_i,+),(L_{j,k,i},-),(T_i,+)]$: each of the four connecting paths ends at a crossing edge in $I_i$ or $T_i$ that lies on $P$; no edge;
    \item $[(I_i,-),(L_{j,k,i},+),(T_i,-)]$: the first and the last connecting path give the edges $uI_i^{\circ}$ and $T_i^{\circ}v$; the other two end at the crossing edge in $L_{j,k,i}$, which lies on $P$.
  \end{itemize}
  If $\tau(x_i)=\mathsf{false}$, all signs are reversed: the first two threads exchange their roles, and so do the last two.
  In both cases the synchronization threads of $x_i$ give only edges $uI_i^{\circ}$ and $T_i^{\circ}v$, and no edge at the unused path of a literal chain link.

  \item \emph{Using a clause thread.}
  The clause thread of $C_k = (\ell_1 \vee \ell_2 \vee \ell_3)$ passes through the literal links $L_{1,k}$, $L_{2,k}$, $L_{3,k}$.
  Since $\tau$ satisfies $C_k$, at least one of its crossing edges lies on $P$, so the clause thread itself is interrupted before it reaches $v$.
  At a crossing edge that is not on $P$, a path can leave the clause thread and enter the unused path of the literal link $L_{j,k,i}$.
  Apart from the clause thread, this unused path is crossed only by a 3-synchronization thread, and the two connecting paths of this thread that start there lead to crossing edges in $I_i$ and $T_i$ that lie on $P$, i.e., to dead ends.

  For the edges of $\Gamma$, recall from the proof of Lemma~\ref{lem:clause-threads} that the crossing edge of the clause thread in $L_{j,k}$ lies on $P$ if and only if $\tau$ satisfies $\ell_j$.
  Hence the four connecting paths of the clause thread give the edge $uL_{1,k}^{\circ}$ if $\ell_1$ is false, the edge $L_{j,k}^{\circ}L_{j+1,k}^{\circ}$ if $\ell_j$ and $\ell_{j+1}$ are false ($j=1,2$), and the edge $L_{3,k}^{\circ}v$ if $\ell_3$ is false, all under $\tau$.
\end{enumerate}

\emph{Conclusion.}
By the three cases, the node $I_i^{\circ}$ is adjacent in $\Gamma$ only to $u$, the node $T_i^{\circ}$ only to $v$, and the unused paths of literal chain links only through edges given by clause threads.
No edge of $\Gamma$ joins $u$ and $v$, because every thread has at least two crossing edges.
Suppose that $\Gamma$ contains a $u$--$v$ walk, and let $W$ be a shortest one; it is a path with at least one internal node.
An internal node of $W$ has two distinct neighbors on $W$, so it is neither $I_i^{\circ}$ nor $T_i^{\circ}$; hence all internal nodes are unused paths of literal chain links.
An edge between two such nodes joins $L_{j,k}^{\circ}$ and $L_{j+1,k}^{\circ}$ for the same clause $C_k$, the node $u$ is adjacent only to nodes $L_{1,k}^{\circ}$ among them, and $v$ only to nodes $L_{3,k}^{\circ}$.
Therefore $W = u\,L_{1,k}^{\circ}\,L_{2,k}^{\circ}\,L_{3,k}^{\circ}\,v$ for some clause $C_k$, and by Case~3 all three literals of $C_k$ are false under $\tau$, a contradiction.
Hence $\Gamma$ contains no $u$--$v$ walk, so all possible continuations from $u$ are blocked in $G\setminus P$.
There is no $u$--$v$ path after removing $P$; hence $P$ is separating.
\end{proof}

\begin{theorem}
\label{thm:ssp-np-complete}
The \SSP{} problem is NP-complete.
\end{theorem}

\begin{proof}
The reduction is Algorithm~\ref{alg:reduction}; we prove its correctness, bound its running time and show that \SSP{} belongs to NP.

\paragraph{Correctness.}
The constructed instance of \SSP{} is equivalent to the original \ThreeSAT{} instance:
\begin{enumerate}[(a)]
    \item If there exists a separating shortest path $P$ in $G$, then a satisfying assignment for the original \ThreeSAT{} instance can be extracted from $P$.

    Since $P$ is a shortest $u$--$v$ path, it is a chain path by Lemma~\ref{lem:chain-paths}, and since it is separating, it hits every thread.
    By the synchronization property (Lemma~\ref{lem:synchronization}), $P$ is consistent, so it determines a truth assignment $\tau$.
    Since $P$ also hits every clause thread, $\tau$ satisfies at least one literal in every clause by Lemma~\ref{lem:clause-threads}.
    Thus, $\tau$ is a satisfying assignment for the original formula.

    \item If there exists a satisfying assignment for the \ThreeSAT{} instance, then a separating shortest path in $G$ can be constructed from it.

    Given a satisfying assignment $\tau$, let $P$ be the consistent chain path that corresponds to $\tau$: in all chain links of the block of $x_i$, it traverses the positive path if $\tau(x_i)=\mathsf{true}$ and the negative path if $\tau(x_i)=\mathsf{false}$.
    The path $P$ is a shortest $u$--$v$ path (Lemma~\ref{lem:chain-paths}), and it hits every thread (Lemmas~\ref{lem:synchronization} and~\ref{lem:clause-threads}).
    By Lemma~\ref{lem:separating}, $P$ is separating, i.e., $G\setminus P$ contains no $u$--$v$ path.
\end{enumerate}

\paragraph{Running time of the reduction.}
The reduction runs in polynomial time.
The chain has $2n+3m$ chain links, and the construction creates $2n+7m$ threads, each of which has at most three crossing edges and at most four connecting paths.
Before the calibration, the chain therefore has $O(n+m)$ edges.
Balancing the two paths of a chain link adds at most two edges for each of its crossing edges, so $\Lambda = O(n+m)$.
Finally, each of the $O(n+m)$ connecting paths is subdivided until it has $\Lambda$ edges, which adds $O((n+m)^2)$ edges in total.
Hence, the size of $G$ and the running time of Algorithm~\ref{alg:reduction} are polynomial in the input size.

\paragraph{Membership in NP.}
A candidate path $P$ is verified in polynomial time: check that $P$ is a shortest $u$--$v$ path, remove its edges from $G$, and test by breadth-first search whether a $u$--$v$ path remains.
Hence \SSP{} belongs to NP.
\end{proof}

\subsection{From Separating Shortest Path to Min Cut-Path}\label{sec:ssp-to-mcp}
The second step reduces \SSP{} to \MinCutPath, which we state in decision form.

\begin{problem}[\MinCutPath{} (Decision Version)]
Let $G=(V,E)$ be an undirected graph and $u,v \in V$ distinct vertices; cut-paths are defined in Definition~\ref{def:cut-path}.
The decision version of the \MinCutPath{} problem is defined as follows:

\begin{center}
\begin{tabular}{p{0.12\linewidth}p{0.8\linewidth}}
\textbf{Input:} & A graph $G=(V,E)$, vertices $u,v \in V$, and an integer $k$. \\
\textbf{Question:} & Does there exist a cut-path $F \subseteq E$ between $u$ and $v$ with $|F|\leq k$?
\end{tabular}
\end{center}
\end{problem}

\begin{theorem}
\label{thm:mcp-np-complete}
The decision version of the \MinCutPath{} problem is NP-complete.
\end{theorem}

\begin{proof}
We reduce \SSP{} to the decision version of \MinCutPath{} and use Theorem~\ref{thm:ssp-np-complete}.

\paragraph{Reduction.}
Let $G$ be an instance of \SSP{} with distinguished vertices $u,v$.
We may assume that $u$ and $v$ lie in the same connected component of $G$, since otherwise the answer is negative.
We define $G' := G$ and set the threshold $k := d_G(u,v)$, the length of a shortest $u$--$v$ path in $G$.
We then ask whether $G'$ admits a cut-path of size at most $k$.
The threshold $k$ is computed by breadth-first search, so the reduction runs in polynomial time.

\paragraph{Correctness.}
Suppose that $G$ admits a separating shortest path $P$.
Then $|P| = d_G(u,v) = k$, and $P$ itself is a cut-path of size $k$, since (i) it is a $u$--$v$ path and (ii) its removal disconnects $u$ and $v$, so $P$ is also a $u$--$v$ cut.
Conversely, suppose that $G'$ has a cut-path $F$ with $|F| \leq k$.
Then $F$ contains a $u$--$v$ path $P$, which has at least $d_G(u,v) = k$ edges; hence $F = P$ and $P$ is a shortest $u$--$v$ path in $G$.
Since $F$ also contains a $u$--$v$ cut, $G \setminus P$ contains no $u$--$v$ path.
Hence $P$ is a separating shortest path, and the two instances are equivalent.

\paragraph{Membership in NP.}
A candidate cut-path $F$ can be verified in polynomial time:
check that $|F| \leq k$, that $F$ contains a $u$--$v$ path, and that $G\setminus F$ contains no $u$--$v$ path (e.g., by breadth-first search).
All conditions can be tested in polynomial time.

\paragraph{Conclusion.}
By Theorem~\ref{thm:ssp-np-complete}, \SSP{} is NP-complete, and it reduces to \MinCutPath{} in polynomial time.
Since \MinCutPath{} is in NP (by polynomial verification), it follows that \MinCutPath{} is NP-complete.
\end{proof}

The decision version is NP-complete, and hence the optimization version is NP-hard.
